% ---------------------------------------------------------------------
%  TÓM TẮT
% ---------------------------------------------------------------------
\begin{abstract}

Dị ứng thực phẩm là vấn đề cần được quan tâm vì việc tiếp xúc hoặc tiêu thụ các thành phần gây dị ứng có thể ảnh hưởng đến sức khỏe người dùng. Tuy nhiên, việc kiểm tra nhãn mác và bảng thành phần trên bao bì thực phẩm bằng phương pháp thủ công có thể mất thời gian, đặc biệt khi thông tin được trình bày bằng chữ nhỏ hoặc sử dụng các tên thành phần khó nhận biết. Xuất phát từ thực tế đó, đồ án xây dựng hệ thống đọc nhãn mác và cảnh báo thành phần dị ứng thực phẩm nhằm hỗ trợ người dùng nhận diện thông tin trên bao bì và phát hiện các thành phần có khả năng gây dị ứng.

Hệ thống cho phép người dùng cung cấp hình ảnh nhãn mác thực phẩm, sau đó sử dụng công nghệ nhận dạng ký tự quang học (OCR) để trích xuất văn bản từ hình ảnh. Nội dung thu được được xử lý và phân tích để xác định các thành phần có liên quan đến dị ứng, kết hợp với dữ liệu về chất gây dị ứng và cơ chế hỗ trợ phân tích của trí tuệ nhân tạo. Dựa trên kết quả nhận diện, hệ thống hiển thị thông tin thành phần và đưa ra cảnh báo tương ứng để người dùng có thêm cơ sở tham khảo trước khi lựa chọn thực phẩm. Đồ án hướng đến việc xây dựng một công cụ hỗ trợ tra cứu thuận tiện, góp phần nâng cao nhận thức về an toàn thực phẩm và giảm bớt công sức kiểm tra nhãn mác thủ công. Kết quả của đề tài là cơ sở để tiếp tục hoàn thiện hệ thống về khả năng nhận dạng, độ chính xác khi phân tích thành phần và phạm vi dữ liệu dị ứng được hỗ trợ.

\vspace{0.8em}
\noindent\textbf{Từ khoá:} OCR, nhận dạng ký tự, dị ứng thực phẩm, phân tích thành phần, trí tuệ nhân tạo, cảnh báo dị ứng.


\end{abstract}
